%=====================================================================
\chapter{The Machine Learning Project}\label{ch:project}
%=====================================================================
\locator{6}{Part VI --- Responsible Practice}

\begin{outcomes}
By the end of this chapter you will be able to:
\begin{tight}
  \item \textbf{Scope} a machine learning project: turn a vague aim into a task,
        performance measure, data source, baseline and success criterion.
  \item \textbf{Plan} the project in milestones and \textbf{write} a one-page
        proposal.
  \item \textbf{Run} experiments reproducibly: one fixed split, cross-validated
        comparisons, an experiment log, seeds and version control.
  \item \textbf{Report} results honestly, with a baseline, uncertainty, per-group
        performance and a model card.
  \item \textbf{Present} and \textbf{defend} the project's choices.
\end{tight}
\end{outcomes}

\begin{prereq}
Everything. In particular \chref{introduction} (T, P and E), \chref{workflow} (the
workflow and leakage), \chref{generalization} (model selection),
\chref{evaluation} (metrics and uncertainty) and \chref{ethics} (the checklist and the
model card). The project delivers HEC 2023 CLO-5 --- ``develop a reasonable size project
using suitable machine learning technique'' --- and the HEC 2025 outcome ``design and
develop end-to-end ML solutions for real-world datasets''.
\end{prereq}

\begin{keyterms}
problem framing $\bullet$ success criterion $\bullet$ proposal $\bullet$ milestone
$\bullet$ fixed test set $\bullet$ experiment log $\bullet$ reproducibility
$\bullet$ random seed $\bullet$ version control $\bullet$ ablation $\bullet$ error
analysis $\bullet$ model card $\bullet$ technical report $\bullet$ presentation
\end{keyterms}

\section{What the Project Is For}

Each chapter has taught one piece: one algorithm, one metric, one pitfall. Real work
needs all of them at once, in the right order, on data nobody has cleaned in advance,
for a question that matters to someone. The project is where that happens. It is
worth 15\% of the course (see the Course Specification), is done in pairs, runs from
Week~11 to Week~16, and is judged less on the final accuracy than on whether every
number in it can be trusted.

\section{Choosing and Scoping a Problem}

A good project problem has four properties: a \textbf{decision} someone would make
with the prediction; \textbf{data} you can legally obtain, with labels, preferably at
least a thousand examples; a \textbf{baseline} that is easy to compute; and a
\textbf{scope} that fits five weeks. Sources of data:
\begin{tight}
  \item the issue trackers of open-source projects (the Apache Software Foundation's
        public Jira, GitHub issues), which record tickets, their text, their history
        and how long they took;
  \item public software-engineering datasets --- effort-estimation and defect data
        such as those of the PROMISE repository --- and the general collections of the
        UCI Machine Learning Repository, OpenML and Kaggle;
  \item data from your own employer's project tracker --- only with written
        permission, and never with personal identifiers (\chref{ethics}).
\end{tight}

\begin{conceptbox}[Project Ideas]
\begin{tight}
  \item \textbf{Resolution time.} Predict how long a ticket will take to resolve
        from its text and metadata (\chref{linreg}, with the word features of
        \chref{bayes}), and compare with the team's own estimates.
  \item \textbf{Ticket routing.} Send each ticket to the right component or team
        from its text; report macro-averaged metrics, because the rare queues matter
        (\chref{evaluation}).
  \item \textbf{Effort estimation.} Predict story hours from story points and task
        features (\chref{linreg}, \chref{knn}), and test whether the model beats
        ``hours per point'' from last quarter.
  \item \textbf{Projects at risk.} Classify projects likely to overrun from their
        early metrics --- the worked example below.
  \item \textbf{Defect prediction.} Flag code changes likely to introduce a defect;
        the classes are imbalanced, so report the precision--recall curve.
  \item \textbf{Duplicate reports.} Find earlier tickets that describe the same
        fault (cosine similarity and nearest neighbours, \chref{knn}).
\end{tight}
\end{conceptbox}

\begin{worked}[From a Vague Idea to a Project]
\textbf{Idea:} ``Use machine learning to deliver more projects on time.''

\smallskip
\textbf{Decision.} Who acts on the prediction? The project office, which can review
about 15 projects in depth each quarter.

\textbf{Task $T$.} One month into each project, rank the projects that started this
quarter by their risk of overrunning their deadline by more than 20\%.

\textbf{Experience $E$.} Four years of completed projects (about 900) with their
month-one metrics --- change requests, staffing, velocity, open defects --- and whether
they overran, exported from the company's project tracker with written permission and
without personal identifiers.

\textbf{Performance $P$.} Recall among the top 15 ranked projects (how many eventual
overruns the office would review), with precision reported beside it.

\textbf{Baseline.} Rank by the number of change requests alone --- the rule the office
already uses.

\textbf{Success criterion.} Beat the change-request rule's recall at 15 by at least
three projects on the held-out year, with no project type (web, mobile, data) whose
recall is below 60\% of the others'.

\smallskip
Every later choice --- features, model, threshold, fairness check --- now has
something to be judged against. A project without that paragraph cannot fail, and so
cannot succeed either.
\end{worked}

\section{The Plan}

\dsfig{%
\begin{tikzpicture}[font=\scriptsize,
  bar/.style={rounded corners=2pt, fill=#1!55, draw=#1, minimum height=0.46cm},
  ms/.style={diamond, fill=accent, inner sep=1.8pt}]
\foreach \w in {11,...,16}{
  \pgfmathsetmacro{\xx}{(\w-11)*2.0}
  \node[font=\scriptsize\bfseries, text=primary] at (\xx+1.0,0.55) {Week \w};
  \draw[gray!30] (\xx,0.3) -- (\xx,-5.2);}
\draw[gray!30] (12.0,0.3) -- (12.0,-5.2);
\foreach \y/\lab/\s/\e/\c in {
   0/{Proposal and data}/0/2.0/primary,
   1/{Baseline and pipeline}/1.0/4.0/secondary,
   2/{Experiments and tuning}/2.0/7.0/greenhl,
   3/{Error analysis and fairness check}/5.0/8.5/purplehl,
   4/{Test once; report and model card}/7.5/10.5/goldhl,
   5/{Presentation}/10.0/12.0/tealhl}{
  \node[anchor=east, font=\scriptsize] at (-0.2,-\y*0.8-0.1) {\lab};
  \node[bar=\c, anchor=west, minimum width={(\e-\s)*1cm}] at (\s,-\y*0.8-0.1) {};}
\node[ms] at (2.0,-0.1) {};
\node[font=\tiny, accent, anchor=west] at (2.15,-0.1) {proposal due};
\node[ms] at (7.0,-1.7) {};
\node[font=\tiny, accent, anchor=west] at (7.15,-1.7) {experiments frozen};
\node[ms] at (10.5,-3.3) {};
\node[font=\tiny, accent, anchor=west] at (10.62,-2.95) {report due};
\end{tikzpicture}}%
{The project across the last six weeks of the semester. The test set is opened only once,
after experiments are frozen in Week~14.}{project-plan}

\begin{conceptbox}[The One-Page Proposal (due end of Week 11)]
\begin{tightnum}
  \item \textbf{Problem and decision:} who will use the prediction, for what.
  \item \textbf{$T$, $P$, $E$}, as in the worked example.
  \item \textbf{Data:} source, size, label, permission, known problems.
  \item \textbf{Baseline} and \textbf{success criterion}.
  \item \textbf{Planned models} (at least one from each of two different chapters).
  \item \textbf{Risks and ethics:} who could be harmed; which fairness check.
\end{tightnum}
\end{conceptbox}

\section{Running Experiments Honestly}

\begin{definitionbox}[Rules for Reproducible Experiments]
\begin{tight}
  \item \textbf{Split once, first,} and store the test indices; the test set is not
        opened until the experiments are frozen.
  \item \textbf{Every model is a pipeline} (\chref{workflow}) evaluated by the same
        cross-validation folds.
  \item \textbf{Log every run}: date, model, settings, CV mean and standard deviation,
        and one line on why it was run. The log, not memory, is the record.
  \item \textbf{Fix the random seeds} and record the library versions.
  \item \textbf{Use version control} (Git): commit the code with each logged run, so
        any number in the report can be regenerated.
\end{tight}
\end{definitionbox}

\begin{worked}[Reading an Experiment Log]
The lab's template writes this log for the projects-at-risk example, whose metric is
recall on the overrun class:

\smallskip
\begin{center}
\begin{tabular}{@{}lrr@{}}
\toprule
\textbf{Model} & \textbf{CV recall (mean)} & \textbf{(sd)} \\
\midrule
baseline (majority class) & 0.000 & 0.000 \\
logistic regression & 0.799 & 0.047 \\
random forest & 0.804 & 0.061 \\
\bottomrule
\end{tabular}
\end{center}

\smallskip
\textbf{The baseline scores zero}: it predicts ``on time'' for every project, so it
never catches an overrun --- the accuracy paradox of \chref{evaluation} expressed in the
metric that matters. \textbf{The two models} --- logistic regression (\chref{logreg})
and a random forest (\chref{ensembles}) --- \textbf{are indistinguishable}: 0.799 and
0.804 differ by far less than one fold-to-fold standard deviation. Choose the simpler,
more explainable one --- logistic regression, whose odds ratios the project office can
read --- tune it (the template finds $C = 0.1$ with balanced class weights, which lifts
cross-validated recall to 0.881), and only then open the test set, once: recall 0.857,
which is 48 of the 56 overruns caught.
\end{worked}

\textbf{Error analysis} is the step most projects skip and the one that teaches most:
look at the examples the model gets wrong, and ask what they have in common. An
\textbf{ablation} --- removing one feature or one step and measuring the change --- shows
what each part of the pipeline contributes. Both belong in the report.

\section{The Report and the Presentation}

\rowcolors{2}{white}{lightbg}
\begin{longtable}{@{}L{3.6cm}L{11.4cm}@{}}
\toprule
\rowcolor{primary!15}
\textbf{Section} & \textbf{Contents} \\
\midrule
\endhead
1. Problem & The decision, $T$, $P$, $E$, success criterion (from the proposal, updated) \\
2. Data & Source, permission, size, label, cleaning, the split; a table of features \\
3. Method & Baseline; models compared; the pipeline; how hyperparameters were chosen \\
4. Results & The experiment log; the chosen model's test performance \emph{with a
confidence interval} (\chref{evaluation}); comparison with the baseline and success
criterion \\
5. Analysis & Error analysis; ablation; per-group performance and the fairness check \\
6. Limitations & What the model must not be used for; what would change the conclusion \\
Appendix & The model card (\chref{ethics}); how to rerun everything \\
\bottomrule
\end{longtable}

The presentation is ten minutes and five slides: the decision and why it matters; the
data; what was compared and how; the result against the baseline, with its uncertainty;
and one limitation. Every member of the pair must be able to answer any question on any
part.

\section{Assessment}

\rowcolors{2}{white}{lightbg}
\begin{longtable}{@{}L{4.4cm}L{1.3cm}L{7.4cm}L{1.6cm}@{}}
\toprule
\rowcolor{primary!15}
\textbf{Criterion} & \textbf{Marks} & \textbf{Full marks require} & \textbf{CLO} \\
\midrule
\endhead
Framing and proposal & 10 & A decision, $T$/$P$/$E$, baseline and success criterion
stated before any modelling & 1, 6 \\
Data and pipeline & 15 & Legal data; leak-free pipeline; one fixed split & 6 \\
Experiments & 20 & At least two model families; CV comparison; logged; reproducible from
the repository & 2--4, 6 \\
Evaluation & 20 & Right metric; baseline; confidence interval; test set used once & 5 \\
Analysis and ethics & 15 & Error analysis; per-group results; model card & 5, 7 \\
Report & 10 & Clear, complete, honest about limitations & 6 \\
Presentation and defence & 10 & Both members answer confidently & 6 \\
\midrule
\textbf{Total} & \textbf{100} & Weighted to 15\% of the course & \\
\bottomrule
\end{longtable}

\begin{pitfall}
\begin{tight}
  \item \textbf{Choosing the dataset before the question.} Start from a decision;
        find data for it.
  \item \textbf{Opening the test set to ``check progress''.} It is then a validation
        set, and the final number is optimistic.
  \item \textbf{No baseline.} A 92\% model means nothing until the baseline is known.
  \item \textbf{Reporting only the best run.} Report the log, including what did not
        work.
  \item \textbf{A notebook that runs only in the order it was written.} Restart and run
        everything before submitting; better, turn it into a script.
\end{tight}
\end{pitfall}

%---------------------------------------------------------------------
\begin{lab}[A Project Template, End to End]
Copy this script as the skeleton of your project and replace the data and the
candidate models. It writes an experiment log, the saved model and a model card into
a folder called \texttt{project\_output}.
\begin{lstlisting}[language=Python]
"""Project skeleton: from raw data to a saved, documented model."""
import csv
import datetime
import json
import os

import joblib
import sklearn
from sklearn.datasets import make_classification
from sklearn.dummy import DummyClassifier
from sklearn.ensemble import RandomForestClassifier
from sklearn.linear_model import LogisticRegression
from sklearn.metrics import classification_report, recall_score
from sklearn.model_selection import (GridSearchCV, StratifiedKFold,
                                     cross_val_score, train_test_split)
from sklearn.pipeline import make_pipeline
from sklearn.preprocessing import StandardScaler

SEED = 42                                   # one seed, used everywhere
OUT = "project_output"
os.makedirs(OUT, exist_ok=True)

# --- 1. data, and a split made once and never changed --------------------
# a stand-in for your tracker export: 900 projects, 8 month-one metrics
X, y = make_classification(n_samples=900, n_features=8, n_informative=5,
                           n_redundant=1, weights=[0.75], flip_y=0.02,
                           random_state=5)  # y: 1 = overran, must not miss
X_tr, X_te, y_tr, y_te = train_test_split(
    X, y, test_size=0.25, stratify=y, random_state=SEED)
cv = StratifiedKFold(5, shuffle=True, random_state=SEED)
metric = "recall"                           # chosen in the proposal

# --- 2. baseline, then candidates, all by cross-validation ---------------
candidates = {
    "baseline": DummyClassifier(strategy="most_frequent"),
    "logistic": make_pipeline(StandardScaler(), LogisticRegression()),
    "forest": RandomForestClassifier(n_estimators=300, random_state=SEED),
}
with open(os.path.join(OUT, "experiments.csv"), "w", newline="") as fh:
    log = csv.writer(fh)
    log.writerow(["model", f"cv_{metric}_mean", f"cv_{metric}_std"])
    for name, model in candidates.items():
        s = cross_val_score(model, X_tr, y_tr, cv=cv, scoring=metric)
        log.writerow([name, round(s.mean(), 4), round(s.std(), 4)])
        print(f"{name:9s} CV {metric} {s.mean():.3f} +/- {s.std():.3f}")

# --- 3. tune the chosen family -------------------------------------------
params = {"logisticregression__C": [0.01, 0.1, 1, 10],
          "logisticregression__class_weight": [None, "balanced"]}
grid = GridSearchCV(candidates["logistic"], params, cv=cv,
                    scoring=metric).fit(X_tr, y_tr)
print("chosen settings:", grid.best_params_)

# --- 4. the test set, once -----------------------------------------------
pred = grid.predict(X_te)
test_recall = recall_score(y_te, pred)
print(classification_report(y_te, pred,
                            target_names=["on time", "overran"], digits=3))

# --- 5. save the model and write its model card --------------------------
joblib.dump(grid.best_estimator_, os.path.join(OUT, "model.joblib"))
card = {
    "model": "logistic regression, standardised features",
    "intended use": "rank projects for a project-office review; never to "
                    "judge a team or a manager",
    "data": f"{len(y)} past projects, 8 month-one metrics",
    "split": f"75/25 stratified, seed {SEED}; 5-fold CV for selection",
    "settings": grid.best_params_,
    f"test {metric}": round(float(test_recall), 3),
    "limitations": "one company's projects; re-validate after any change "
                   "to how projects are run",
    "date": datetime.date.today().isoformat(),
    "library": f"scikit-learn {sklearn.__version__}",
}
with open(os.path.join(OUT, "model_card.json"), "w") as fh:
    json.dump(card, fh, indent=2, default=str)
print("saved:", sorted(os.listdir(OUT)))
\end{lstlisting}
The template reproduces the experiment log of the worked example. Before adapting it,
change \texttt{metric} to \texttt{"precision"} and rerun: does the chosen model change?
\end{lab}

%---------------------------------------------------------------------
\begin{checkpoint}
\begin{tightnum}
  \item Write a success criterion for the late-task warning of \chref{introduction}
        that a project manager could check.
  \item Two models score 0.912 $\pm$ 0.030 and 0.905 $\pm$ 0.028 in cross-validation.
        Which do you choose, and on what grounds?
  \item Name three things that must be recorded for an experiment to be reproducible.
\end{tightnum}
\end{checkpoint}

%---------------------------------------------------------------------
\begin{chaptersummary}
\begin{tight}
  \item Start from a decision, and write $T$, $P$, $E$, a baseline and a success
        criterion before modelling.
  \item Split once; compare pipelines by cross-validation on fixed folds; log every
        run; fix seeds; commit code.
  \item Choose between near-equal models on simplicity; open the test set once;
        report the result with its uncertainty and beside the baseline.
  \item Error analysis, ablations and per-group results turn a score into
        understanding.
  \item The report, the model card and the defence are part of the work, not
        decoration.
\end{tight}
\end{chaptersummary}

\begin{reviewq}
  \item \textbf{(MCQ)} The test set should be used: (a)~after every experiment
        (b)~once, after experiments are frozen (c)~to choose hyperparameters
        (d)~never.
  \item \textbf{(MCQ)} An ablation study: (a)~removes components to measure their
        contribution (b)~adds more data (c)~tunes the learning rate (d)~cleans labels.
  \item \textbf{(MCQ)} The most important part of a project proposal is: (a)~the
        choice of algorithm (b)~the decision, metric and success criterion (c)~the
        number of features (d)~the programming language.
  \item \textbf{(Short)} Why should an experiment log include runs that did not
        improve the model?
  \item \textbf{(Short)} What is error analysis, and what can it reveal that a single
        metric cannot?
  \item \textbf{(Short)} Explain why the worked example chose logistic regression over
        the random forest.
  \item \textbf{(Applied)} Write the one-page proposal for an IT project-management
        problem of your choice, using the template of this chapter.
  \item \textbf{(Applied)} For the worked example's test result (48 of 56 overruns
        caught), compute a 95\% confidence interval for the recall and discuss what
        it means for the claim ``the model catches 86\% of overruns''.
\end{reviewq}
